\begin{abstract}
The likelihood of encountering in-training failures rises substantially with larger Deep Learning (DL) training workloads, leading to lost work and resource wastage. Such failures are typically offset by checkpointing, which comes at the cost of storage and network bandwidth overhead. State-of-the-art approaches involve lossy model compression mechanisms, which induce a tradeoff between the resulting model quality and compression ratio.
We make a key enabling observation that the sensitivity of model weights to compression varies during training, and different weights benefit from different quantization levels, ranging from retaining full precision to pruning. 
We propose (1) a non-uniform quantization scheme that leverages this variation, 
(2) an efficient search mechanism that dynamically finds the best quantization configurations, and
(3) a quantization-aware delta compression mechanism that rearranges weights to minimize checkpoint differences and thereby improving compression.

We instantiate these contributions in \sys, an in-training checkpoint compression system for DL workloads. Our experiments show that \sys consistently achieves a better tradeoff between accuracy and compression ratio compared to prior works, enabling a compression ratio up to 39x and withstanding up to 10 restores with negligible accuracy impact in fault-tolerant training. \sys achieves at least an order of magnitude reduction in checkpoint size for failure recovery and transfer learning without any loss of accuracy.



\end{abstract}
